%Special Problem manuscript (Chapters I-III) based on icsthesissample.tex
%Template author: JCEPlaras
%Follows APA citation

\documentclass{icsthesis}
\usepackage[T1]{fontenc}
\usepackage{graphicx}
\usepackage{amsmath}
\usepackage{array}
\usepackage{url}
\usepackage{nameref}

%\usepackage{showframe} %for debugging the margins

%Cross-reference to a section by its title (headings are unnumbered)
\newcommand{\secref}[1]{\textit{\nameref{#1}}}

%icsthesis.cls turns hyphenation off, so allow looser word spacing instead of letting long words
%run past the right margin
\setlength{\emergencystretch}{5em}

%Replace the values here, various commands depend on this variables
\renewcommand{\TITLE}{ORDINAL DEEP LEARNING FOR AUTOMATED FALL ARMYWORM DAMAGE SEVERITY SCORING IN
FAWCHECK: A CNN-BASED APPROACH TO THE PRASANNA SCALE}
\renewcommand{\AUTHOR}{GIOVANN JOSH N. APOLINAR}
\renewcommand{\DEGREE}{BACHELOR OF SCIENCE IN COMPUTER SCIENCE}
\renewcommand{\MONTH}{NOVEMBER}
\renewcommand{\YEAR}{2026}
\newcommand{\ADVISER}{ARIAN J. JACILDO}

\begin{document}

	\begin{frontmatter}
		%create title
		\maketitle

		%TOC
		\maketableofcontents

		%Resync the class's page counter, since the TOC can run past one page
		\setcounter{pagecounter}{\value{page}}\addtocounter{pagecounter}{-1}

		%LOT
		\makelistoftables

		%ABSTRACT (not needed for the proposal; uncomment to restore)
		\setcounter{pagecounter}{\value{page}}\addtocounter{pagecounter}{-1}
		\begin{abstractwithpageno}
		\\
		\textbf{\AUTHOR}, University of the Philippines Los Ba\~{n}os, \MONTH\ \YEAR.
		\textbf{\TITLE}
		\\Adviser: \ADVISER\\
		(abstract text here)
		\end{abstractwithpageno}

	\end{frontmatter}

	%BODY OF THESIS HERE (Use up to 3 levels only, (sec -> subsec -> subsubsec)
	\begin{mainmatter}

		%=====================================================================
		\section{INTRODUCTION}\label{sec:intro}

			\subsection{Background of the Study}\label{sec:background}

				The fall armyworm (FAW), \textit{Spodoptera frugiperda} (J.~E. Smith), is among the
				most invasive pests that bring serious economic loss to Philippine maize production.
				The Department of Agriculture recorded infestations affecting more than 5,000
				farmers across approximately 3,200 hectares by mid-2024, as reported in the FAWCheck
				v1 study \citep{casasola2026}. By mid-August 2024, the pest had destroyed an
				estimated 734 metric tons of corn worth about PHP~57.03 million. These losses
				contributed to a national deficit, with second-quarter output falling 20.35\%. Farm
				grain prices, meanwhile, remained nearly flat while retail corn prices rose 5.75\%,
				leaving growers selling less grain of lower quality with no price gain to offset the
				loss. Because of the FAW's ability to spread quickly, the time between detection and
				intervention can often mean the difference between a manageable infestation and a
				lost crop. Effective integrated pest management necessitates the ability to
				determine not only whether a plant is infested, but also the extent of damage the
				plant has sustained, as varying degrees are important in evaluating germplasm
				resistance.

				Typically, the screening process uses a visual rating scale as the basis for
				evaluation. The scale used in this study is the one to nine foliar damage scale of
				\citet{prasanna2018}, which is modified from the whorl-stage screening scales of
				\citet{davis1992}. A field technician assesses the plant, matches what is visible
				against a set of ordered descriptors, and assigns an integer score.

				Thus, the main motivation of this study lies in two characteristics of this rating
				scale. First, the scale is ordinal and not nominal. A score of 5 denotes a
				significantly greater degree of damage than a score of 4. A model that treats the
				nine categorical scores as unordered labels discards this ranking structure. With
				the use of multi-class cross-entropy, predicting 3 when the true target is 6 incurs
				the same loss as predicting 5 when the true target is 6. Essentially, such a loss
				function is blind to order, which is fundamental in this problem.

				Second, the distance between scores carries meaning in the screening process.
				\citet{prasanna2018} present the scale for assessing foliar damage in maize
				germplasm entries, and group its scores into resistance-response bands: score 1
				highly resistant, scores 2 and 3 resistant, scores 4 and 5 partially resistant,
				scores 6 and 7 susceptible, and scores 8 and 9 highly susceptible (see
				Table~\ref{tab:prasanna}). In assessing the resistance of maize, a confusion between
				6 and 7 leaves an entry in the susceptible category, while a confusion between 5 and
				7 moves the resistance label to a different categorical response, from partially
				resistant to susceptible. Therefore, errors on this scale are not interchangeable. A
				model that is blind to these differences is likely to produce inefficient results
				regardless of how accurate its metrics appear.

				Moreover, FAWCheck v1 was built as a field monitoring tool for agricultural
				technicians and researchers \citep{casasola2026}. In this use, the score describes
				how badly a crop is damaged. The most expensive error is scoring a susceptible or
				highly susceptible plant (scores 6 through 9) below 6, because it understates damage
				where damage is greatest. This study therefore treats scores 6 through 9 as
				\textit{severe damage}, starting where the susceptible category of
				\citet{prasanna2018} begins. Over-scoring also has a cost, because it overstates
				damage. The study measures both whether models keep plants in the correct response
				category, in either direction, and whether they detect severe damage (see
				\secref{sec:metrics}).

			\subsection{Statement of the Problem}\label{sec:problem}

				FAWCheck v1 is an existing application that digitizes this rating workflow. It is a
				progressive web application which uses a convolutional neural network (CNN) to score
				an image of maize, along with a rule-based scoring system that feeds on a symptom
				questionnaire. The outputs of the two components are fused into a final severity
				estimate. It is a working system with measured performance, and the present study
				builds on this foundation rather than displacing it. The CNN was reported at 50.00\%
				exact accuracy, 0.4105 macro F1, 76.47\% within-one-level accuracy, and 0.8088 mean
				absolute error (MAE). Moreover, the fused system reached 83.13\% exact agreement
				with expert scores, 0.27 MAE, and 0.7881 weighted Cohen's kappa
				\citep{casasola2026}.

				However, three problems limit what can be built on those results. First, the v1 CNN
				treats severity scoring as multi-class classification with sparse categorical
				cross-entropy \citep{casasola2026}. It cannot represent the inherent order of the
				severity scale, and its loss penalizes the categories with similar magnitude without
				accounting for the real distance among target scores. Nonetheless, the errors were
				already largely adjacent to the true target, as 76.47\% of its predictions fell
				within one level of the reference \citep{casasola2026}. This implies that the errors
				did not arise from random noise but are structured around the true label. Thus, the
				question is not whether ordinal structure matters, but whether modeling it reduces
				the more distant errors that remain.

				Second, the v1 results cannot serve as a baseline. The CNN metrics were computed on
				a small evaluation set of 68 images. Applying the Wilson score method to the
				published counts gives a 95\% confidence interval (CI) of [0.3844, 0.6156] for the
				CNN accuracy, an interval width of 0.2312. These findings do not point to a defect
				in the model, but to a limit of what a sample of 68 images can establish. At that
				sample size, the 95\% CI spans roughly 38\% to 62\%, wide enough to contain both a
				poor classifier and a good one. Consequently, any later model that scores within
				this range cannot be said to have improved on the 50\% baseline, since the
				difference would be indistinguishable from sampling noise at this evaluation size.

				Third, the v1 study correctly attributes the weak performance of the CNN to class
				imbalance, limited samples at some severity levels, variation in field image
				quality, and visual overlap between neighboring categories \citep{casasola2026}.
				Indeed, the dataset is severely skewed toward the lower-end scoring categories, but
				this is not the only factor that contributes to the weak performance of the model. A
				review of the file-level structure of the train/validation/test partition, conducted
				for this study, shows that the partition is non-representative (see
				\secref{sec:discard}). Because validation results also guided model selection in v1
				\citep{casasola2026}, that decision process inherited the same bias from the
				non-representative partition. This is a separate defect from class imbalance, and
				any before-and-after comparison run on this partition measures artifacts of the
				partition rather than genuine differences in model performance.

			\subsection{Objectives}\label{sec:objectives}

				To determine whether ordinal regression improves automated FAW foliar damage
				severity scoring on the Prasanna et al.\ one to nine scale relative to standard
				multi-class classification, under an evaluation protocol capable of supporting that
				comparison.

				Specifically, the researcher would like to do the following:

				\begin{enumerate}
					\item To compare the rank-consistent CORAL \citep{cao2020} and CORN
						\citep{shi2023corn} and the regression-based THOR \citep{fuchs2022} against
						a softmax multi-class baseline on identical cross-validation folds, with
						identical imbalance handling \citep{kang2020} applied to all four models,
						and determine whether ordinal regression reduces severity-scoring error.
					\item To determine whether the three ordinal methods differ from one another on
						this task, and whether the published expectations---CORN advantaged on mean
						absolute error and CORAL on outlier-sensitive error \citep{shi2023corn},
						measured here as macro-averaged mean squared error (see
						\secref{sec:evaluation}), and THOR advantaged over CORAL on mean absolute
						error \citep{fuchs2022}---hold under severe class imbalance at nine ordinal
						levels.
					\item To measure the agreement between the author's and an expert's independent
						scores on every newly labeled image, to locate the model's agreement with
						the expert relative to it, and to check the consistency of the inherited
						labels across near-duplicate images.
					\item To rebuild the FAWCheck evaluation protocol on a stratified,
						leakage-controlled foundation, and determine how much of the change in
						measured performance is attributable to the protocol rather than to the
						model, by training the same v1 recipe on the same images under the inherited
						partition and under the rebuilt protocol.
				\end{enumerate}

				The third objective sets the models beside a human rater: each model's agreement with
				the expert is reported next to the author's, on the same images (see
				\secref{sec:humanref}). This is a reference point, not a ceiling (see
				\secref{sec:scope}).

			\subsection{Significance of the Study}\label{sec:significance}

				For the Institute of Plant Breeding, the Entomology Laboratory Division, and the
				users of FAWCheck, this study produces a foliar damage severity model whose
				performance is reported with quantified uncertainty: confidence intervals on
				accuracy, macro-MAE, and weighted kappa, taken from repeated stratified
				cross-validation with corrected variance and from bootstrap resampling by
				near-duplicate cluster. The v1 study reported its performance as a point estimate,
				an accuracy of 50.00\% on 68 evaluation images, with no interval. This quantified
				uncertainty is the precondition for two decisions. First, it indicates whether the
				image-based score can be trusted in the field. Second, it indicates whether future
				dataset growth actually improves the model, which can be established if a future
				result clears the interval of the earlier one.

				For agricultural computer vision, the study applies rank-consistent ordinal
				regression to an expert-rated agronomic damage scale. As discussed in
				\secref{sec:cnnsev}, existing crop severity studies mostly use a small number of
				coarse bins evaluated with plain accuracy. Meanwhile, ordinal methods and
				agreement-based metrics are established practice in medical image grading.

				Methodologically, the study contributes a reusable evaluation protocol, described in
				\secref{sec:protocol}, that can serve subsequent FAWCheck model iterations and other
				rating-scale tasks in the Laboratory's work. The protocol extends the v1 study,
				whose own recommendations called for expanded, better-balanced data
				\citep{casasola2026}.

			\subsection{Scope and Limitations of the Study}\label{sec:scope}

				The study only covers the image-based severity model of FAWCheck. It does not
				address cob or ear damage, yield loss, the rule-based heuristic scoring, the fusion
				of image-based and rule-based scores, or the integration of the resulting model into FAWCheck's deployed inference service, including its export to the ONNX format. Whether fusion improves on the image model
				alone is recorded as an open question but is not tested. The v1 comparison of its CNN-only and fused-system figures \citep{casasola2026} cannot settle it: the two figures come from different datasets, class mixes, and reference labels, and the rule engine's symptom answers mirror the descriptors behind the reference score. Settling it requires image-only, rule-only, and fused arms on the same held-out cases. The scope only covers
				foliar damage on the Prasanna et al.\ one to nine severity scale
				\citep{prasanna2018}.

				Moreover, the model comparison is limited to a softmax baseline and three ordinal
				methods: CORAL and CORN from the rank-consistent binary-decomposition family, and
				THOR from the regression-based family. It does not cover the whole ordinal
				regression family. Soft ordinal labels are excluded, and the regression-based family
				is represented by THOR alone, not also by plain regression with tuned cut-points
				(see \secref{sec:whythor}). Replacing MobileNetV3Small with a stronger self-supervised backbone, such as the distilled DINOv3 ViT-S/16 \citep{simeoni2025}, is left to future work.

				Furthermore, the prediction target is an image-based Prasanna rating: a score
				assigned by a rater looking at a photograph, not at the plant. The descriptors count
				affected leaves and refer to the whorl and furl, which a single photograph may show
				only in part, so both the labels and the agreement measured in this study describe
				rating photographs, not scoring plants in the field. The scale was also designed for
				maize germplasm entries, typically assessed under controlled infestation
				\citep{prasanna2018,davis1992}, whereas much of the imagery here comes from farmer
				fields and public datasets. FAWCheck applies the scale to field monitoring, an
				extension of its stated purpose, so the response categories are read as levels of
				damage, not as judgments of a plant's resistance. Because the scale is the derived
				form identified in the review of rating scales (see \secref{sec:scales}), results
				reported here are not strictly comparable to published screening results reported
				against the original Davis or Williams scales.

				Whether the ordinal comparison is tested as a confirmatory hypothesis depends on a
				sensitivity check made before any ordinal model is trained (see
				\secref{sec:sensitivity}). If the comparison cannot detect the smallest effect of
				interest even after the pre-registered responses, it is reported as exploratory.

				Score 9 is represented by four images in the inherited dataset. The expansion target
				requires it to reach at least 15; if it does not, it is merged with score 8 for
				reporting under a pre-registered rule and is out of scope for quantitative per-level
				claims (see \secref{sec:merge}). Even at 15 images, per-level claims for score 9
				carry wide intervals.

				In addition, model performance depends on the representativeness of the training
				data. The inherited dataset mixes field photographs with publicly sourced images,
				and the candidate expansion pool is drawn from a single region and season in
				Tanzania, captured largely on non-phone cameras. Generalization to Philippine field
				conditions is therefore measured, not assumed, on a Philippine external test set of
				locally collected phone photographs (see \secref{sec:holdout}). At a minimum of 150
				images, a proportion near 50\% on that test set carries a 95\% confidence interval
				about 16 percentage points wide, against 23 for v1's 68-image evaluation.

				Finally, no human ceiling is measured. Agreement is measured between the author and
				one expert, which partly reflects how well the author has learned the rubric and is
				therefore a reference point rather than a bound on achievable performance (see
				\secref{sec:annotation}). A ceiling would need two experts scoring the same images
				independently, which this study does not do and the literature does not report for
				this scale (see \secref{sec:observer}). The inherited 2,084 labels, whose provenance is
				undocumented \citep{casasola2026}, are not re-scored. Their quality is assessed only
				through their consistency across near-duplicate images, which can reveal
				inconsistency but cannot establish accuracy.

		%=====================================================================
		\section{REVIEW OF RELATED LITERATURE}\label{sec:rrl}

			This chapter reviews the rating scale and its reliability, prior image-based approaches
			to plant damage and disease severity, ordinal regression methods for deep networks and
			their use in medical grading, and the evaluation methodology under which model
			comparisons on imbalanced ordinal data can be defended. It closes by stating the gap the
			present study addresses.

			\subsection{Fall Armyworm Damage Rating Scales}\label{sec:scales}

				The 1--9 foliar damage scale used by FAWCheck derives from the visual rating scales
				developed by \citet{davis1992} for screening whorl-stage maize for FAW resistance,
				and reaches this study through \citet{prasanna2018}. The descriptors are cumulative
				in character: lower scores describe shot-hole feeding on a few leaves, middle scores
				introduce elongated lesions of increasing extent, and upper scores describe lesions
				across most leaves accompanied by large irregular holes, terminating at near-total
				destruction of the whorl and furl.

				\citet{toepfer2021} review the family of FAW leaf damage rating scales and draw a
				distinction that matters for this study: Davis' 0--9 scale assesses the whorl and
				furl from a top view, whereas Williams' 0--9 scale assesses the whole plant. The two
				are frequently integrated in applied work, and \citeauthor{toepfer2021} observe that
				the scales have been repeatedly modified, and at times misapplied or miscited
				\citep{toepfer2021}. They propose simplified,
				instar-independent variants on the grounds that the descriptive distinctions in the
				original scales are difficult to apply consistently in the field.

				\textit{Scale identification.} The operative scale for this study is the one
				published as Table 3 of \citet{prasanna2018} and condensed in
				Table~\ref{tab:prasanna}. Its source line reads ``modified from Davis and Williams
				(1992)'' \citep{davis1992}. It is therefore neither Davis' whorl-and-furl top-view
				scale nor Williams' whole-plant scale in its original form, but a derived scale
				drawing on both, and it is named as such here rather than claimed to correspond
				strictly to either predecessor.

				\begin{table}[htbp]
				\vspace{4ex}
				\centering
					\caption{Prasanna scale for foliar damage by Fall Armyworm, with descriptors}
					\label{tab:prasanna}
					\singlespacing\footnotesize
					\begin{tabular}{cp{9.2cm}l}
					\hline
					\hline
					Score & Damage symptoms & Response$^{\star}$ \\ \hline
					1 & No visible leaf-feeding damage & Highly resistant \\[3pt]
					2 & Few pinholes on one or two older leaves & Resistant \\[3pt]
					3 & Several shot-holes on a few leaves (fewer than 5) and small circular holes &
					Resistant \\[3pt]
					4 & Shot-holes on several leaves (6--8), or small lesions and pinholes, plus a
					few small elongated lesions (up to 1.3~cm) on whorl and furl leaves & Partially
					resistant \\[3pt]
					5 & Elongated lesions (over 2.5~cm) on 8--10 leaves, plus a few small to
					mid-sized irregular holes eaten from the whorl or furl & Partially resistant
					\\[3pt]
					6 & Several large elongated lesions and/or large irregular holes on several
					whorl and furl leaves & Susceptible \\[3pt]
					7 & Many elongated lesions of all sizes on several whorl and furl leaves, plus
					several large irregular holes & Susceptible \\[3pt]
					8 & Many elongated lesions of all sizes on most whorl and furl leaves, plus many
					mid- to large-sized irregular holes & Highly susceptible \\[3pt]
					9 & Whorl and furl leaves almost totally destroyed; plant dying from extensive
					leaf damage & Highly susceptible \\ \hline\hline
					\end{tabular}

					\vspace{1ex}
					\begin{minipage}{0.85\textwidth}
					\singlespacing\footnotesize
					Adapted from Table 3 of \citet{prasanna2018}, whose source line reads ``modified
					from Davis and Williams (1992).'' $^{\star}$Response labels are read in order
					from the printed column, whose entries are offset from their rows. This reading
					assigns score 1 to highly resistant, consistent with other descriptions of the
					scale; it is to be confirmed against a clean copy of the guide.
					\end{minipage}
				\vspace{4ex}
				\end{table}

				The published table supports reading it as admitting both assessment framings.
				First, its descriptors mix whole-plant leaf counting with whorl-specific detail:
				scores 2 through 5 are defined by how many leaves carry injury (``1--2 older
				leaves,'' ``a few leaves (\textless 5 leaves),'' ``several leaves (6--8 leaves),''
				``8--10 leaves''), while scores 4 through 9 add the condition of the whorl and furl
				leaves. Second, the nine reference photographs accompanying it (Figure 5 of
				\citealp{prasanna2018}) show both framings, some filling the frame with one plant's
				upper leaves and others showing whole plants in a stand. An image is accordingly
				taken to satisfy this scale's assessment conditions if it shows the extent of leaf
				injury across the plant, the condition of the whorl and furl, or both, in either
				framing. This is the sense in which the capture protocol's requirement of one
				primary maize plant per photograph, with the whorl and upper leaves in frame, is
				consistent with the scale.

			\subsection{Observer Variability on the Rating Scale}\label{sec:observer}

				\citet{toepfer2021} document observer bias on these scales qualitatively, including
				difficulty in the mid-to-upper range, but report no kappa, intraclass correlation,
				or inter-rater correlation coefficient. A partial quantitative estimate does exist.
				In proposing a simplified TNAU 1--5 scale, \citet{srinivasan2022} report that when
				the same subjects scored the same plots, the mean standard deviation across raters
				ranged from 0.48 to 0.80 under the Davis 1--9 scale, against 0.19 to 0.40 under the
				1--5 alternative. The authors attribute the difference to overlapping descriptions
				between adjacent Davis scores---noting in particular that damage on one-half of
				leaves (score 7) and on two-thirds of leaves (score 8) is difficult to distinguish
				reliably---and observe that reproducibility declines as the number of scale points
				increases \citep{srinivasan2022}.

				Consequently, the claim that the literature supplies no quantitative estimate of
				rater variability on this scale is too strong; a dispersion figure exists, even if
				an agreement coefficient does not. However, the figure that exists is not the figure
				this study needs. A standard deviation across raters scoring plots is not a
				quadratic weighted kappa between raters scoring individual images, and the imagery,
				raters, and conditions differ. The TNAU result is therefore best used as a prior
				expectation against which this study's own measurement is positioned, not as a
				substitute for making it. Its finding that mid-to-upper-range distinctions are least
				reliable also predicts where author--expert disagreement in this study should
				concentrate, which is a testable expectation rather than an assumption.

				A terminological caution follows from the same source: it distinguishes the original
				Davis 1--9 scale from a modified version it attributes to Ni et al.\ (as cited in
				\citealp{srinivasan2022}). Together with the Davis/Williams distinction above, the
				phrase ``the 1--9 scale'' is ambiguous across at least three referents, which is why
				the review of rating scales above names the particular derived scale this study
				uses. It also means that published results reported against ``the Davis scale'' or
				``the Williams scale'' are not strictly comparable to results reported here, a
				limitation recorded in the \secref{sec:scope}.

			\subsection{CNN Approaches to Plant Disease and Pest Severity}\label{sec:cnnsev}

				Deep learning for plant disease has been dominated by identification rather than
				grading. \citet{mohanty2016} and \citet{ferentinos2018} established that CNNs
				classify disease \textit{type} from leaf imagery with high accuracy, and subsequent
				reviews document the breadth of that work \citep{lu2021,boulent2019}. Severity
				assessment is a much smaller literature: a survey of CNN-based plant disease
				severity assessment identifies only sixteen studies and notes that severity work
				remains sparse relative to identification \citep{shi2023severity}.

				The severity studies that exist share a structural pattern relevant here. They
				discretize severity into a small number of coarse bins---four stages for maydis leaf
				blight of maize \citep{haque2022}, four to six stages for tomato bacterial spot and
				mosaic virus \citep{astani2026}, and five bands defined by the proportion of
				diseased leaves in work covered by the survey \citep{shi2023severity}---and they
				evaluate with plain accuracy and F1. The bins are typically derived from measured
				lesion or leaf area rather than assigned by an expert against a descriptor rubric.
				The survey reports that the studies it reviews are evaluated with accuracy,
				precision, recall, mean average precision, and F1, and it describes none that uses
				an ordinal-aware loss or an ordinal error measure such as MAE or quadratic weighted
				kappa \citep{shi2023severity}; the abstracts of \citet{haque2022} and
				\citet{astani2026} likewise report accuracy-based measures only.

				Furthermore, a four-bin severity task derived from lesion area has fewer decision
				boundaries than a nine-point expert-rated task, those boundaries are defined by a
				continuous measurable quantity rather than by a rubric, and its ground truth is not
				subject to inter-rater disagreement in the same way. Results from that literature do
				not transfer to the present task, and the present task's difficulty should not be
				judged against its accuracy figures.

			\subsection{Deep Learning for Fall Armyworm}\label{sec:fawdl}

				Work specific to FAW is, with one exception noted below, binary or
				detection-oriented rather than graded. \citet{sena2003} established an early
				image-based approach, segmenting maize plants and classifying image blocks as
				damaged or non-damaged, with 94.72\% of 720 images correctly classified.
				\citet{ishengoma2021} applied VGG16, VGG19, InceptionV3, and MobileNetV2 to UAV
				imagery for infected-versus-uninfected classification, with a follow-up hybrid CNN
				for faster detection \citep{ishengoma2022}. \citet{zhang2022} address pest
				identification rather than damage, distinguishing \textit{S.~frugiperda} adults and
				larvae among 36 maize field pests. \citet{sandhya2025} fuse RGB and thermal imagery
				for infested-versus-healthy classification, and \citet{prabha2021} deploy a
				MobileNetV2-based CNN in a mobile application that distinguishes healthy from
				infested plants, reporting 87\% validation accuracy. Reported accuracies across
				these studies range from 87\% \citep{prabha2021} to 100\%
				\citep{ishengoma2021}---figures that reflect the comparative ease of a two-class or
				identification problem rather than transferable performance on severity grading.

				The nearest approach to severity is that of \citet{dzurume2025}, who classify FAW
				infestation severity in Bangladeshi maize fields into three levels (none, moderate,
				severe) using temporal Sentinel-2 imagery and a random forest classifier, reporting
				overall accuracy between 74\% and 84\%. This is severity work, but at field scale,
				from satellite multispectral data, with three bins and a non-neural classifier; it
				addresses neither plant-level assessment, the nine-point scale, nor ordinal
				structure.

				No published work was located that applies CNNs to ordinal grading of FAW foliar
				damage on the Davis/Prasanna 1--9 scale. FAWCheck v1 applies a CNN to the scale
				itself, but as nine-class classification \citep{casasola2026}.

			\subsection{Ordinal Regression for Deep Networks}\label{sec:ordinal}

				\citet{niu2016} introduced the binary-decomposition approach for CNNs---converting a
				$K$-class ordinal target into $K-1$ binary ``is the rank greater than $k$?''
				subtasks---and demonstrated it on age estimation. The approach suffers a structural
				defect: the independently trained binary classifiers can produce inconsistent
				predictions, estimating a higher probability that the rank exceeds 2 than that it
				exceeds 1.

				\citet{cao2020} resolve this with the Consistent Rank Logits (CORAL) framework, in
				which the $K-1$ binary tasks share a single weight vector and differ only in their
				bias terms. They prove that this weight sharing forces the predicted cumulative
				probabilities into non-increasing order, guaranteeing rank consistency by
				construction rather than by post-hoc correction, and confirm it empirically: zero
				rank inconsistencies on every dataset and seed, against consistently nonzero counts
				for the unconstrained predecessor \citep{cao2020}.

				CORAL's structure, a single weight vector shared by all thresholds with a separate
				bias for each, is the neural counterpart of the proportional-odds model of
				\citet{mccullagh1980}, in which one set of coefficients is shared across all
				cumulative splits of the scale and only the intercepts differ.

				\citet{shi2023corn} argue that CORAL's guarantee is bought at a cost---the shared
				weight vector restricts the expressiveness of the output layer---and propose CORN,
				which achieves rank consistency through conditional training subsets and the chain
				rule for conditional probabilities, removing the weight-sharing constraint. They
				report a substantial gain over CORAL. Only in supplementary experiments on two
				text-review datasets did CORAL attain the lower root mean squared error while CORN
				attained the lower MAE, which the authors attribute to RMSE penalizing large errors
				more heavily; on their image datasets, CORN is lower on both \citep{shi2023corn}.

				A third family abandons binary decomposition and regresses a single scalar score,
				assigning a class according to where that score falls among class boundaries. The
				threshold-based ranking loss (THOR) of \citet{fuchs2022} fixes those boundaries in
				advance at integer values and trains the regression with a pairwise ranking loss, so
				that inputs are projected onto the segment of their class. Its authors report the
				best MAE among the ordinal methods they compared, which included CORAL
				\citep{fuchs2022}. No official implementation of THOR is published. On five-grade diabetic retinopathy with a ResNet18 backbone, THOR reached MAE 0.709 against CORAL's 0.965 \citep{fuchs2022}, the most direct published evidence against a rank-consistent head. Plain regression with cut-points tuned after training is the simpler member of the family, but \citet{fuchs2022} include no plain-regression baseline, so their evidence does not rank THOR against it; THOR's fixed boundaries do remove the post-training cut-point step, which would otherwise have to be kept free of evaluation data. Soft ordinal labels \citep{diaz2019} encode the ordinal prior in the label distribution rather than the output layer, a separate line of work.

				\citet{gutierrez2016} provide the broader taxonomy---na\"ive approaches, ordinal
				binary decompositions, and threshold models---across sixteen algorithms and
				forty-one datasets. The survey predates the CORAL/CORN line and contains no deep
				learning methods, so it is cited here for its taxonomy and methodological standard
				rather than for algorithmic guidance.

				An important qualification applies to both CORAL and CORN as evidence for the
				present study. Both were validated on face-age estimation at 33 to 55 ranks. CORN's
				primary datasets were rebalanced to near-equal class counts before training, and
				neither paper reports a macro-averaged or otherwise imbalance-aware metric
				\citep{cao2020,shi2023corn}. Neither tests small $K$, and neither tests a
				distribution with the tail mass of this study's data. THOR's evidence shares the
				limitation: its five-class experiments did test small $K$, but the diabetic
				retinopathy and aesthetics datasets were subsampled to equal class sizes before
				training \citep{fuchs2022}. The rank-consistency guarantee is a proof and transfers
				regardless; the reported magnitude of predictive improvement does not. No method is therefore an expected winner on a nine-level, heavily skewed scale. The one \textit{a priori} consideration is that CORAL's weight-sharing restriction, the expressiveness cost identified by \citet{shi2023corn}, applies to a backbone fine-tuned with each method's own loss, an untested configuration in which that trade-off is most likely to favor CORN.

				None of these methods addresses class imbalance in training, and in this study
				imbalance is handled separately from ordinal structure. \citet{kang2020} show that a
				network's learning can be divided into two parts, representation and classifier.
				They find that imbalance need not harm representation learning, that representations
				learned on the natural, unrebalanced distribution generalize best, and that strong
				performance on rare classes can be recovered by adjusting only the classifier. Their
				classifier retraining (cRT) method keeps the trained representation fixed,
				re-initializes the classifier, and retrains only the classifier on class-balanced
				samples \citep{kang2020}. Because it acts only on the final layer, it applies in the
				same form to any output layer, which is what this study requires.

			\subsection{Ordinal Methods in Medical Grading}\label{sec:medical}

				Where agricultural severity work has largely not adopted ordinal methods, medical
				image grading has. Grading against clinical severity scales is common in that field,
				and quadratic weighted kappa is its conventional reporting metric even for models
				trained as ordinary classifiers: a cross-entropy-trained network for diabetic
				retinopathy reports quadratic weighted kappa of 0.888 \citep{farag2022}, and a
				cascaded CNN ensemble for five-step AREDS macular degeneration staging reports
				quadratic kappa of 0.870 \citep{xie2023}. Ordinal formulations proper are also
				established: ordinal regularization for joint diabetic retinopathy and lung nodule
				grading \citep{tang2023}, rank-consistent CORAL heads in current retinopathy systems
				\citep{braganca2026}, and deep probabilistic ordinal regression for diabetic
				retinopathy and prostate cancer grading \citep{toledo2022}.

				The one located instance of an ordinal severity head in plant imaging is instructive
				as a contrast case rather than a precedent. \citet{bohra2026} propose a multi-task
				vision transformer for rose leaf disease that includes an ordinal severity head
				alongside classification and uncertainty estimation. In their limitations, the
				authors state that the severity labels ``originate from disease categories and not
				from expert-approved measurement systems'' \citep{bohra2026}, although a later
				passage describes the annotations as expert-provided; the paper is not internally
				consistent on this point. Its ablation results, however, settle the practical
				question: the full model reports severity MAE of $0.000 \pm 0.000$ with Spearman
				$\rho$ of 1.000 \citep{bohra2026}, indicating that severity is a deterministic
				function of the class label. The ordinal task is therefore solved trivially and
				carries no information about ordinal modeling under genuine label uncertainty.

			\subsection{Dataset Bias and Leakage in Plant Imaging}\label{sec:leakage}

				Reported accuracy in plant disease imaging is unusually vulnerable to dataset
				artifacts. \citet{noyan2022} demonstrates the mechanism on PlantVillage: a random
				forest trained on only eight background pixels per image reached 49.0\% accuracy
				across 38 classes, against a 2.6\% chance baseline. \citet{mohanty2016} had earlier
				reported that their PlantVillage-trained network fell from approximately 99\% to
				approximately 31\% accuracy when tested on images from other online sources, a
				result \citet{noyan2022} also cites. Background and capture-condition artifacts
				correlate with class, and a model can learn the correlation instead of the
				pathology.

				This has two direct consequences for the present study. First, the v1 dataset mixes
				approximately 1,500 field and client-provided photographs with approximately 584
				publicly sourced images, and per-image provenance is not currently recorded
				\citep{casasola2026}; if source correlates with severity, a model can learn a
				dataset signature. Second, the v1 dataset retained multiple images of the same plant
				where available to increase variation \citep{casasola2026}, which is precisely the
				pattern that leaks across randomly assigned splits. Both require explicit control in the evaluation protocol. The size of the leakage effect is measurable: in CIFAR, replacing test images that had near-duplicates in the training set lowered accuracy by 9--14\% relative \citep{barz2020}.

Shift between training and deployment imagery can also be anticipated in training. Training a new output layer on frozen features before fine-tuning keeps a randomly initialized layer from distorting pretrained features, with reported gains of about 1\% in distribution and 10\% out of distribution \citep{kumar2022}. Augmenting with the corruptions of the deployment device helps only for those corruptions: networks trained on one distortion become robust to it but generalize poorly to others \citep{geirhos2018}. JPEG compression and motion blur are among the common corruptions cataloged by \citet{hendrycks2019}.

			\subsection{Evaluation and Comparison Methods for Imbalanced Ordinal
			Data}\label{sec:evalmethods}

				\citet{baccianella2009} propose macro-averaged variants of MAE, mean squared error,
				and zero-one error---computing the error separately per true class and then
				averaging across classes---and prove these coincide with the conventional
				micro-averaged forms when class cardinalities are equal, while resisting trivial
				classifiers on imbalanced data. Their diagnostic argument is the substantive one:
				under micro-averaged MAE on skewed data, the error-minimizing constant prediction
				need not even be the majority class, so a system that never predicts rare classes
				can outscore one that attempts them.

				The same argument applies to quadratic weighted kappa, which is sensitive to
				marginal distributions in the same way. With scores 6 through 9, the susceptible and
				highly susceptible categories, comprising 14.0\% of this dataset, kappa could reward
				a model that declines to predict severe damage---the specific failure this
				application cannot tolerate.

				Macro-averaging has a countervailing cost the source paper does not address: giving
				every class equal weight regardless of support means a class with few test items can
				move the average substantially. \citet{baccianella2009} validate on one dataset
				family, one learner, and one right-skewed distribution whose rarest class, though
				3.9\% of the set, still numbers in the hundreds. Its behavior at the per-class test
				counts arising in this study is undocumented, which is a reason to measure its
				stability empirically rather than assume it.

				For comparing models, \citet{dietterich1998} established that paired statistical
				tests applied across cross-validation folds are anti-conservative, because folds
				sharing training data are not independent. \citet{nadeau2003} provide the corrected
				resampled \textit{t}-test, which inflates the variance estimate to account for
				training-set overlap, and show that a variance estimate that ignores this overlap
				is too small, making spurious significant differences between algorithms likely. \citet{bouckaert2004} found this test replicable, and most
				replicable when applied to ten-times-repeated ten-fold cross-validation.

				\citet{gutierrez2016} supply the methodological benchmark this study is held
				against: nested cross-validation to prevent hyperparameter leakage into the test
				estimate, twenty to thirty random splits per dataset, and Wilcoxon signed-rank tests
				with Bonferroni correction. The survey also reports its own external-validity
				limitation---its discretized-regression datasets are binned into equal-frequency
				classes and so exhibit no class imbalance, and performance patterns diverge between
				the discretized and genuinely ordinal halves of its benchmark \citep{gutierrez2016}.
				That is the disclosure the balanced-evaluation choices elsewhere in the ordinal
				literature needed and did not make.

			\subsection{Reporting Performance Against a Human Ceiling}\label{sec:ceiling}

				Where model predictions are compared against expert labels that themselves carry
				disagreement, agreement among experts bounds what the model can achieve, and
				reporting it alongside model performance is established practice in clinical
				imaging. \citet{schaap2022} provide a clear template: in automating Psoriasis Area
				and Severity Index scoring, they report the intraclass correlation between CNN and
				real-life physician scores, report the agreement among five PASI-trained physicians
				scoring the same images (ICC 0.706--0.793), and compare the CNN's agreement with the
				real-life scores against the physicians' image-based agreement with them.

				None of the ordinal methodology papers reviewed above provides an analog; all treat
				their ground truth as uncontested. The present study cannot measure an expert
				ceiling, because it does not have two experts score the same images independently.
				It adopts a weaker form of the pattern of \citet{schaap2022}: model agreement with
				an expert is set beside the agreement between the same expert and the author, a
				rater trained on the rubric, on the same images.

			\subsection{Gaps Addressed by the Present Study}\label{sec:gaps}

				The literature supports four observations that jointly motivate this study.

				First, ordinal structure in severity scales is real and consequential, and
				rank-consistent formulations to exploit it exist and are proven
				\citep{cao2020,shi2023corn}. Those formulations, however, have been validated at
				large $K$ on balanced or micro-averaged data, never at nine levels under severe
				skew.

				Second, agricultural severity work has not adopted them. It uses coarse bins derived
				from measurable area or from class identity, evaluated with plain accuracy
				\citep{shi2023severity,haque2022,astani2026,bohra2026}. FAW-specific work is binary
				or detection-oriented
				\citep{sena2003,ishengoma2021,ishengoma2022,zhang2022,%
				sandhya2025,prabha2021,dzurume2025},
				and the one CNN applied to the Prasanna scale treats it as nominal
				\citep{casasola2026}.

				Third, ordinal regression is used in medical grading against expert-rated scales
				with measurable rater disagreement \citep{tang2023,braganca2026,toledo2022}, a field
				where agreement-based reporting is conventional
				\citep{farag2022,xie2023,schaap2022}. No located work applies rank-consistent
				ordinal regression to an expert-rated, agronomically standardized damage scale in
				crop imaging. The Prasanna 1--9 scale sits in that gap.

				Fourth, the evaluation standards under which an improvement claim can be
				defended---stratified leakage-controlled splits, macro-averaged metrics, corrected
				significance testing, and a measured human ceiling---are established in the
				methodological and clinical literature
				\citep{gutierrez2016,noyan2022,baccianella2009,dietterich1998,%
				nadeau2003,bouckaert2004,schaap2022}
				but are not yet standard in agricultural computer vision.

				This study applies the fourth set of standards to the first set of methods, on a
				task where neither has previously been brought to bear, with one substitution: an
				author--expert agreement reference in place of a measured human ceiling.

		%=====================================================================
		\section{MATERIALS AND METHODS}\label{sec:methods}

			This chapter describes the research design, data sources, annotation procedure,
			evaluation protocol, FAWCheck v1 baseline, models, training procedure, evaluation methods, and statistical analysis used to compare
			ordinal and multi-class severity models for FAWCheck. Every rule in this chapter,
			including readouts, the merge rule, the tuning budget, decision rules, and the
			responses to a failed sensitivity check, is fixed before any ordinal model is
			trained, and any deviation is reported. The four models compared receive identical
			folds, backbone, imbalance handling, augmentation, tuning budget, and computing
			environment.

			\subsection{Research Design}\label{sec:design}

				The study follows a comparative experimental design consisting of:

				\begin{enumerate}
					\item \textbf{Audit and rebuilding of the evaluation protocol}
						(see \secref{sec:protocol}).
					\item \textbf{Expansion and labeling of the dataset}, prioritizing the upper
						severity range, and \textbf{collection of a Philippine external test set}
						(see \secref{sec:data}).
					\item \textbf{Measurement of author--expert agreement} and a
						\textbf{label-consistency check} (see \secref{sec:annotation}).
					\item \textbf{Re-measurement of the current FAWCheck recipe}, ported to PyTorch,
						to measure the protocol's effect, establish the baseline, and supply the
						variance for a sensitivity check (see \secref{sec:v1comp} and
						\secref{sec:sensitivity}).
					\item \textbf{Comparison of softmax classification, CORAL, CORN, and THOR} on
						MobileNetV3Small, with corrected significance testing (see
						\secref{sec:models} and \secref{sec:evaluation}).
					\item \textbf{A single final evaluation on the Philippine external test set}
						(see \secref{sec:holdout}).
				\end{enumerate}

			\subsection{Materials}\label{sec:materials}

				The hardware and software used are listed in Table~\ref{tab:materials}. All runs that enter a single
				model comparison are made in the same computing environment, and every run records the
				environment it ran on.

				\begin{table}[htbp]
				\vspace{4ex}
				\centering
					\caption{Hardware and software}
					\label{tab:materials}
					\singlespacing\footnotesize
					\begin{tabular}{p{3.0cm}p{6.4cm}p{4.4cm}}
					\hline
					\hline
					Item & Specification & Role \\ \hline
					Primary workstation & ASUS ROG Strix G614JV: Intel Core i7-13650HX (20 threads), 48~GB RAM,
					NVIDIA GeForce RTX 4060 Laptop GPU (8~GB GDDR6), 2~TB storage, Ubuntu 26.04.1 LTS & All
					training and evaluation, by default \\
					Fallback & NVIDIA T4 GPU (16~GB GDDR6), through Google Colab & Only when the workstation is
					unavailable or a run exceeds its GPU memory \\ \hline
					Python, PyTorch & --- & Every model, including the re-measured v1 baseline \\
					coral-pytorch & Reference implementation by the methods' authors \citep{raschka2022coral} &
					CORAL and CORN layers and losses \\
					THOR & Implemented from its paper; no published code \citep{fuchs2022} & THOR output layer and
					loss \\
					Keras (TensorFlow) & Inference only, on the v1 model file & Arm (a) of the protocol-effect design \\
					Albumentations & \citet{buslaev2020} & Augmentation \\
					scikit-learn & Stratified group $k$-fold splitter & Fold construction \\
					imagehash & 64-bit DCT perceptual hash & Near-duplicate clustering \\ \hline\hline
					\end{tabular}
				\vspace{4ex}
				\end{table}

			\subsection{Data Sources}\label{sec:data}

				\subsubsection{Inherited Labeled Dataset}\label{sec:inherited}

					The primary dataset is the 2,084-image labeled working set assembled for
					FAWCheck v1, each image assigned an integer score from 1 to 9 against the
					Prasanna et al.\ foliar damage descriptors \citep{casasola2026,prasanna2018}.
					Its composition by severity is given in Table~\ref{tab:inherited}. Scores 1
					through 4 account for 75.8\% of the set, scores 7 through 9 for 7.4\%, and score
					9 for four images.

					\begin{table}[htbp]
					\vspace{4ex}
					\centering
						\caption{Severity distribution of the inherited labeled dataset}
						\label{tab:inherited}
						\begin{tabular}{cccl}
						\hline
						\hline
						Score & Images & \% of set & Resistance response$^{\star}$ \\ \hline
						1 & 385 & 18.5\% & Highly resistant \\
						2 & 318 & 15.3\% & Resistant \\
						3 & 487 & 23.4\% & Resistant \\
						4 & 389 & 18.7\% & Partially resistant \\
						5 & 214 & 10.3\% & Partially resistant \\
						6 & 137 & 6.6\% & Susceptible \\
						7 & 98 & 4.7\% & Susceptible \\
						8 & 52 & 2.5\% & Highly susceptible \\
						9 & 4 & 0.2\% & Highly susceptible \\ \hline
						Total & 2,084 & 100\% & \\ \hline\hline
						\end{tabular}

						\vspace{1ex}
						\begin{minipage}{0.85\textwidth}
						\singlespacing\footnotesize
						$^{\star}$Response categories as listed in Table~\ref{tab:prasanna}.
						\end{minipage}
					\vspace{4ex}
					\end{table}

					Its source composition and missing per-image provenance are described under
					\secref{sec:leakage}; provenance is reconstructed before training (see
					\secref{sec:leakctrl}).

				\subsubsection{Expansion Pool}\label{sec:expansion}

					An unlabeled expansion pool is available, comprising the \texttt{maize\_faw}
					subset of the YEESI Lab Dataset \citep{fue2022}---collected in Morogoro,
					Tanzania, in August and September 2022, and released under a CC-BY-4.0
					license---together with images collected locally by the project team. Not all of
					the 726 YEESI images are usable for this study, leaving approximately 579.
					Expansion labeling prioritizes the score 6--9 range, where the inherited dataset
					is thinnest.

					The YEESI imagery differs from the
					deployment context in cultivar, soil, background vegetation, and ambient light,
					and 337 of its approximately 579 usable images were captured on an action camera
					rather than a phone. It is therefore used as training data, not as a validation
					or test population.

					\textit{Expansion target.} Labeling continues until every score holds at least
					100 images in the cross-validation pool, counting inherited and newly labeled
					images together and excluding the Philippine external test set. Where a score
					cannot reach 100 from the available sources, its minimum is 15 images, the
					threshold of the merge rule. On current counts, scores 1 through 6 already
					exceed 100 (Table~\ref{tab:inherited}). Score 7 needs at least 2 more images,
					score 8 at least 48, and score 9 at least 96, or at least 11 to reach the
					minimum.

				\subsubsection{Philippine External Test Set}\label{sec:holdout}

					An external test set of at least 150 images is drawn exclusively from new
					photographs collected locally by the project team on phones, the deployment
					device. None comes from the inherited dataset or the YEESI pool, so the external
					test set is independent of both v1's labeling and the Tanzanian imagery.
					External test images are labeled through the same pipeline as the expansion pool
					(see \secref{sec:labeling}). They are reserved before folds are built, excluded
					from cross-validation and from every model-selection step, and evaluated once,
					after the final models are trained. No external test image may share a
					near-duplicate cluster with any cross-validation image, and a plant photographed
					more than once contributes to only one side.

					The external test set keeps the severity distribution of the photographs as
					collected rather than being balanced across scores, since its purpose is to
					estimate performance on the images FAWCheck will actually receive. It therefore
					supports overall measures and severe-damage recall, with wide intervals at the
					rare scores, but not per-score claims at scores 8 and 9.

			\subsection{Expert Annotation, Agreement, and Label Consistency}\label{sec:annotation}

				\subsubsection{Labeling Procedure for New Images}\label{sec:labeling}

					Every newly labeled image, in the expansion pool and the Philippine external
					test set, is scored independently by two people against the Prasanna et
					al.\ descriptors \citep{prasanna2018}: the author, and an expert from the IPB
					Entomology Laboratory. The expert scores each image before seeing the author's
					score, so that the two scores are independent and free of anchoring. The expert
					then sees the author's score and assigns the final label, which is the label
					used for training and evaluation.

					Each image's three scores are recorded in an expert correction log: the author's
					score, the expert's independent score, and the final score, with the expert's
					reason wherever the final score differs from either. The log distinguishes two
					kinds of change. First, a \textit{correction} is a final score that differs from
					the author's. Second, a \textit{revision} is a final score that differs from the
					expert's own independent score after seeing the author's. Because revisions move
					the final label toward the author's judgment, their frequency is reported as a
					direct measure of how much the author's score influenced the final labels.

					\textit{Photo adequacy.} Before scoring, the expert checks three conditions for
					each new image: the whorl and furl are visible, the image is in focus, and the
					plant fills enough of the frame for the descriptors to be applied. An image that
					fails any of them is recorded as unscorable and excluded from training and
					evaluation. Unscorable images are counted, and their rate in the external test
					set is reported as the share of field photographs FAWCheck could not score. 

				\subsubsection{Author--Expert Agreement}\label{sec:agreement}

					Agreement between the author's and the expert's independent scores is computed
					on all newly labeled images as quadratic weighted kappa, macro-averaged MAE,
					exact agreement, and within-one-level agreement. It is reported separately for
					the expansion images and the external test set, since expansion labeling
					concentrates on scores 6 through 9 while the external test set follows its
					as-collected distribution, and by score group within each.

					This agreement is a human reference point, not a ceiling (see \secref{sec:scope} and \secref{sec:humanref}). Disagreement is
					expected to concentrate in the mid-to-upper range, where rater reproducibility
					on this scale is reported to be worst \citep{srinivasan2022}.

				\subsubsection{Label Consistency Across Near-Duplicates}\label{sec:consistency}

					The inherited 2,084 labels are not re-scored. Their consistency is checked using
					the near-duplicate clusters formed for fold construction, which are formed
					without reference to labels. Images in one cluster are usually photographs of
					the same plant and should carry the same score. For every cluster of two or more
					images, this check reports whether its labels disagree, the share of multi-image
					clusters that do, and the mean absolute difference between paired labels,
					overall and by score. The same check is applied to the newly labeled images.

					The check can reveal inconsistent labeling but cannot establish that
					labels are accurate: a cluster may hold different plants
					photographed alike, or one plant whose damage changed between
					photographs, and near-duplicates were probably labeled together by
					the same person. Discordant clusters are reported, not relabeled.

			\subsection{Evaluation Protocol}\label{sec:protocol}

				\subsubsection{Why the Inherited Partition Is Discarded}\label{sec:discard}

					A file-level audit of the inherited partition, conducted by direct
					count and hashing of the dataset directories rather than from the v1
					study's reported figures, finds it non-representative (see
					Appendix~B). Score 1 is 11.7\% of training but 47.0\% of
					validation; score 6 has no test images and score 9 none in
					validation or test; and because validation results guided model
					selection in v1 \citep{casasola2026}, checkpoints were selected
					against a distribution dominated by score 1. These defects lie in
					how the data were divided, not in the data or their class imbalance,
					so any comparison computed on this partition measures the partition,
					and it is discarded rather than patched.

				\subsubsection{Repeated Stratified Group Cross-Validation}\label{sec:cv}

					Evaluation uses five-fold cross-validation repeated over three random seeds, for
					fifteen train-and-evaluate runs. Folds are constructed with stratification on
					the merged label (see \secref{sec:merge}) and grouping on near-duplicate cluster
					(see \secref{sec:leakctrl}). Training uses the full nine-level score; the merged
					label governs fold assignment only.

					At $n = 2{,}084$, each fold holds approximately 417 images, including
					approximately 20 score-7 and 11 score-8 images. These supports are thin enough
					to require confidence intervals on per-class metrics and substantial enough for
					those intervals to be meaningful.

					No fold's validation data influence training on that fold, including early
					stopping, checkpoint selection, and any threshold tuning.

				\subsubsection{Leakage and Provenance Controls}\label{sec:leakctrl}

					\begin{itemize}
						\item \textit{Near-duplicate grouping.} The v1 dataset retained
						multiple images of the same plant \citep{casasola2026}, which leak across random splits (see \secref{sec:leakage}).
						Near-duplicate cluster is therefore the primary grouping key.
						Clusters are connected components of a 64-bit perceptual hash at
						a label-blind calibrated threshold, checked by an embedding-based review of a random sample (see Appendix~A). Each
						cluster is indivisible in fold assignment.
						\item \textit{Provenance tagging.} Because source and capture artifacts can
							correlate with class \citep{noyan2022}, each image is tagged with its
							source---v1 field, v1 public, YEESI, or local collection---and
							performance is reported by source. If field-photograph performance is materially worse
							than public-image performance, that gap is the deployment-relevant
							estimate.
						\item \textit{Device tagging.} Capture device is recorded, and performance is reported by device. Device is not a grouping key because,
							where device correlates with severity, device grouping makes
							stratification impossible and yields folds that measure the device gap
							rather than model quality.
					\end{itemize}

			\subsection{FAWCheck v1 Baseline and Protocol Effect}\label{sec:v1comp}

				\subsubsection{Baseline Re-Measurement}\label{sec:baseline}

					The current model is re-measured by retraining its recipe, ported to PyTorch
					(see \secref{sec:recipe}), and that figure---not the v1 published figure---is
					the baseline against which improvement is claimed. The port is read out by
					argmax, as v1 was \citep{casasola2026}, with its median readout also reported.
					Porting keeps the baseline in the same framework and computing environment as
					the four models compared (see \secref{sec:materials}), so a difference between
					them cannot reflect the framework. However, the baseline is then a
					reimplementation. Keras and torchvision distribute separately trained ImageNet
					weights for MobileNetV3Small, each expecting its own input scaling, so the
					re-measured figure describes v1's \textit{recipe}, not the v1 model file.

					The v1 published results are reported for context only in
					Table~\ref{tab:v1perf}, with confidence intervals computed for this study from
					the published counts, since the v1 study reports none \citep{casasola2026}. The two sets of figures rest on different evaluation sets, the v1 set over-representing severe damage about twofold (see Appendix~B), and are not directly comparable.

					\begin{table}[htbp]
					\vspace{4ex}
					\centering
						\caption{FAWCheck v1 reported performance with recomputed 95\% confidence
						intervals (Wilson score method)}
						\label{tab:v1perf}
						\begin{tabular}{lccc}
						\hline
						\hline
						Reported measure & Point estimate & 95\% CI & Width \\ \hline
						CNN exact accuracy (34/68) & 50.00\% & [38.4\%, 61.6\%] & 23.1 pp \\
						CNN within $\pm 1$ level (52/68) & 76.47\% & [65.1\%, 85.0\%] & 19.8 pp \\
						System exact agreement (69/83) & 83.13\% & [73.7\%, 89.7\%] & 16.0 pp \\
						System within $\pm 1$ level (75/83) & 90.36\% & [82.1\%, 95.0\%] & 12.9 pp
						\\ \hline\hline
						\end{tabular}
					\vspace{4ex}
					\end{table}

					The v1 CNN macro F1 of 0.4105 is not carried forward as a target:
					its evaluation set held one image at score 9 and two at score 8, too
					few for a stable per-class F1.

				\subsubsection{The v1 Recipe and Its Port}\label{sec:recipe}

					The recipe is read from the trained v1 model file,
					\texttt{fawcheck\_finetuned.keras}, by parsing its saved
					configuration without executing it, and from the v1 manuscript
					\citep{casasola2026}; Table~\ref{tab:recipe} (see
					Appendix~D) lists each setting and its PyTorch counterpart.
					The file does not record five settings: the number of epochs, the
					batch size, the stopping and checkpoint criterion, any class
					weighting, and whether a phase training only the output layer
					preceded fine-tuning. These are filled with stated defaults, which
					are replaced if v1's training code becomes available. Since the purpose is to measure
					v1's recipe as it was, its quirks are reproduced rather than
					corrected, including a contrast augmentation that only ever
					increases contrast. The baseline receives v1's four augmentations
					only; the phone degradations (see \secref{sec:augment}), linear
					probing before fine-tuning, and classifier retraining belong to this
					study's training procedure and are not applied to it.

				\subsubsection{Separating the Protocol from the Model}\label{sec:armdesign}

					Specific Objective 4 asks how much of a change in measured performance is due to
					the evaluation protocol rather than to the model. The port is therefore trained
					and evaluated four ways, and only adjacent pairs are compared:

					\begin{itemize}
						\item \textit{(a) The v1 model file}, run for inference only, on the
							inherited test split (Table~\ref{tab:partition}).
						\item \textit{(b) The port under the inherited partition:} trained on the
							inherited training split, with checkpoints selected on the inherited
							validation split as in v1 \citep{casasola2026}, and evaluated on the
							inherited test split, over three seeds.
						\item \textit{(c) The port under the rebuilt protocol, on the same images:}
							the 2,084 inherited images only, under the five-fold, three-seed
							stratified group cross-validation (see \secref{sec:cv}), evaluated out
							of fold.
						\item \textit{(d) The port under the rebuilt protocol, on the full
							cross-validation pool} after expansion. This is the baseline described
							under \secref{sec:baseline}.
					\end{itemize}

					The difference between (a) and (b) is the port's fidelity: framework, pretrained checkpoint, the defaults of Table~\ref{tab:recipe}, and seed. The difference between (b) and (c) is the protocol's effect, the
					quantity Objective 4 asks for. The recipe, framework, and images are the same in
					both; what differs is how the images are partitioned, whether near-duplicates
					may straddle training and evaluation, and which images the result is computed
					on. The difference between (c) and (d) reflects the added data. Each figure is
					reported with its cluster-level bootstrap interval, under the same metrics as
					the other models. No significance test is applied to the (b)--(c) comparison,
					since (b) rests on a single partition.

					Two limits apply to arms (a) and (b). First, they use v1's partition
					but not v1's evaluation set: the inherited test split holds 146
					images, whereas v1 reported its CNN results on 68
					\citep{casasola2026}. Second, the test split has no images at scores
					6 and 9 (Table~\ref{tab:partition}), so macro-averaged MAE under (a)
					and (b) averages over the other seven scores only.

			\subsection{Models Compared}\label{sec:models}

				\subsubsection{Comparison Set}\label{sec:compset}

					Four models are compared (Table~\ref{tab:models}), all on the current MobileNetV3Small backbone
					\citep{howard2019}, so that differences between output layers are not confounded with architecture
					and the results remain comparable with v1. Softmax multi-class classification is the control and the
					formulation of the current v1 model \citep{casasola2026}. Because CORAL and CORN guarantee that the
					predicted probability that the score exceeds $k$ does not increase with $k$, their threshold count
					equals the median of the predicted distribution over scores. Softmax has no readout fixed by its
					method, so it is read out in two ways from the same trained model: the median, primary, and the
					argmax, the rule v1 used \citep{casasola2026}.

					\begin{table}[htbp]
					\vspace{4ex}
					\centering
						\caption{Models compared, with output layer, loss, and readout}
						\label{tab:models}
						\singlespacing\footnotesize
						\begin{tabular}{p{2.3cm}p{3.6cm}p{3.1cm}p{4.5cm}}
						\hline
						\hline
						Model & Output layer & Loss & Readout to a score from 1 to 9 \\ \hline
						Softmax (control) & 9 logits & Cross-entropy & Median of the predicted distribution (primary);
						argmax, as in v1, also reported \\
						CORAL \citep{cao2020} & 8 threshold logits; one shared weight vector, a bias per threshold &
						CORAL loss & One plus the number of thresholds $k$ with predicted $P(\text{score} > k) > 0.5$ \\
						CORN \citep{shi2023corn} & 8 conditional logits; no weight sharing & CORN conditional loss & As
						CORAL, with $P(\text{score} > k)$ the cumulative products of the conditional outputs \\
						THOR \citep{fuchs2022} & 1 linear unit & Thresholded pairwise ranking loss & Score whose segment
						$[s-1, s]$ contains the output; outputs beyond the outer boundaries go to 1 or 9 \\ \hline\hline
						\end{tabular}
					\vspace{4ex}
					\end{table}

					The median readout is primary because the median of a predicted
					distribution minimizes absolute error, the comparison metric, and it
					puts softmax, CORAL, and CORN on the same statistic; reading softmax
					by argmax alone would credit ordinal training with a gain due only
					to the readout. The argmax result shows, at no extra training cost,
					how much of any difference from v1 is due to the readout alone.

					No ordinal head is treated as the expected winner (see
					\secref{sec:ordinal}). A finding that ordinal training does not
					improve severity scoring by a meaningful margin meets the first
					specific objective, provided it is established by the equivalence
					test (see \secref{sec:hypotheses}) rather than inferred from a
					non-significant difference.

					\textit{THOR implementation.} THOR has no published code and is
					implemented from its paper's definitions \citep{fuchs2022}. The
					paper's boundaries $\{-1, 0, 1, \ldots, K\}$, for classes indexed
					from zero, become $\{0, 1, \ldots, 9\}$ here, so that score $s$
					occupies $[s-1, s]$. The boundaries and the paper's margin $\gamma =
					0.5$ are fixed, not tuned, since the other three models have no
					output-layer hyperparameters. As in the paper, the loss uses pairs
					of images with adjacent scores: each image in a batch is paired with
					a randomly drawn training image one score higher, or one lower for
					score 9, so the anchors keep the natural distribution in stage 1 and
					the class-balanced distribution in classifier retraining. The output
					layer is a single linear unit rather than the paper's two added
					fully connected layers, to keep its capacity comparable with the
					other three models, and it is trained by this study's procedure (see
					\secref{sec:training}), not the paper's optimizer settings.

				\subsubsection{Scope of the Comparison, and Why THOR}\label{sec:whythor}

					THOR represents the regression-based family, for the reasons given
					in \secref{sec:ordinal}. No claim is made about THOR against plain
					regression with tuned cut-points, which its source does not test
					\citep{fuchs2022}. Soft ordinal labels \citep{diaz2019} are
					excluded.



				\subsubsection{Imbalance Handling}\label{sec:imbalance}

					Imbalance is handled by classifier retraining (cRT) \citep{kang2020}, which
					separates the ordinal-structure effect from the imbalance-handling effect (see
					\secref{sec:ordinal}). It is applied in two stages:

					\begin{itemize}
						\item \textit{Stage 1, representation.} Each model's full network is trained
							with that model's own loss on the natural, instance-balanced
							distribution of the training folds, with no rebalancing.
						\item \textit{Stage 2, classifier retraining.} The backbone is frozen, the
							model's output layer is re-initialized, and only that layer is retrained
							on class-balanced samples, restarting the learning rate, for a short
							schedule \citep{kang2020}.
					\end{itemize}

					Here the output layer is the final layer mapping the backbone's
					pooled features to the model's outputs (Table~\ref{tab:models}).

					cRT takes the same form for all four models, whereas a per-threshold
					loss weighting has a natural form only for CORAL and CORN, and it
					keeps the stage-1 comparison free of reweighting. Hard oversampling
					is not used, since duplicating images from a 52-image class with
					known near-duplicates invites memorization.

				\subsubsection{Pre-Registered Merge Rule}\label{sec:merge}

					The reporting rule is fixed before expansion-pool labels are collected:

					\begin{quote}
						\singlespacing
						Any score with fewer than 15 images in the final cross-validation pool is
						merged with the adjacent score toward the middle of the scale, repeating
						until every merged group holds at least 15 images. Merged groups are used
						for fold stratification, for the selection metric, and for headline metrics;
						per-level results are reported in an appendix.
					\end{quote}

					On current counts, only score 9 ($n = 4$) triggers the rule, forming an \{8, 9\} group of 56. If the expansion target (see \secref{sec:expansion}) is met, the merge does not
					trigger and score 9 enters the headline groups on its own; otherwise the rule
					applies as a safeguard. Its effect on the metric is described under
					\secref{sec:selmetric}.

			\subsection{Training Procedure}\label{sec:training}

				\subsubsection{Linear Probing, then Fine-Tuning}\label{sec:lpft}

					Stage 1 of training follows the linear-probe-then-fine-tune (LP-FT) procedure of
					\citet{kumar2022}. First the pretrained backbone is frozen and only the model's
					output layer is trained, until validation macro-averaged MAE stops improving.
					Then the whole network is unfrozen and fine-tuned at a lower learning rate.

					This guards against the shift between Tanzanian action-camera
					training images (see \secref{sec:expansion}) and Philippine phone
					photographs at test (see \secref{sec:leakage}).

					With stage 2 (see \secref{sec:imbalance}), each model is trained in three
					phases: linear probing, fine-tuning, and classifier retraining, for every model. Tuning and phase lengths are described under
					\secref{sec:tuning}.

				\subsubsection{Model Selection}\label{sec:selection}

					Early stopping and checkpoint selection use validation macro-averaged MAE,
					computed on merged groups, with each model's primary readout (see
					\secref{sec:compset}), and drawn only from the inner-validation split of each
					training fold. The argmax readout of softmax is computed from the checkpoint
					selected under its median readout, so that the two softmax results describe one
					trained model.

				\subsubsection{Augmentation}\label{sec:augment}

					The augmentation follows v1 \citep{casasola2026}, extended for phone capture. v1
					applied random horizontal flipping, random rotation, random zoom, and random
					contrast adjustment, to cover variation in leaf orientation, camera angle,
					distance, and lighting \citep{casasola2026}; all four are retained.

					Three degradations typical of phone photographs are added: JPEG
					compression, motion blur, and downscaling. They match the deployment
					device, for the reason given in \secref{sec:leakage}, and the
					phone-only external test set tests whether the robustness was
					achieved.

					The settings are given in Table~\ref{tab:augment}. The geometric settings are
					those recorded in v1's model file (see \secref{sec:recipe}). Contrast is varied
					symmetrically, up and down by 10\%, as v1's ``random contrast'' was evidently
					intended (the v1 baseline keeps v1's own; see \secref{sec:recipe}). The lower
					JPEG quality bound, $q_{\min}$, is the 5th percentile of the JPEG quality
					estimated from the quantization tables of the locally collected phone
					photographs, excluding the external test set, so that the trained range matches
					the deployment device rather than a guess.

					\begin{table}[htbp]
					\vspace{4ex}
					\centering
						\caption{Augmentation settings, applied in the order listed}
						\label{tab:augment}
						\begin{tabular}{lp{7.5cm}c}
						\hline
						\hline
						Transformation & Setting & Probability \\ \hline
						Horizontal flip & --- & 0.5 \\
						Rotation & Uniform within $\pm 36^{\circ}$, reflect padding & 1.0 \\
						Zoom & Scale 0.9--1.1, aspect ratio preserved, reflect padding & 1.0 \\
						Contrast & Factor 0.9--1.1 about each channel's mean & 1.0 \\
						Motion blur & Kernel 3--7 pixels & 0.2 \\
						Downscaling & Factor 0.5--0.9, then restored to $224 \times 224$ & 0.3 \\
						JPEG compression & Quality $q_{\min}$--95 & 0.5 \\ \hline\hline
						\end{tabular}
					\vspace{4ex}
					\end{table}

					JPEG compression is applied last, because compression is the last step of a
					phone's own image pipeline.

				Two augmentations are excluded as not label-preserving, the property
				\citet{shorten2019} call an augmentation's \textit{safety}: vertical
				flips, which v1 also did not use \citep{casasola2026}, since whole-plant
				photographs have a fixed stem-below, leaves-above orientation; and
				channel shuffling and other color-space transformations, since leaf
				color may carry information about damage and plant condition
				\citep{shorten2019}. Contrast adjustment changes contrast without
				altering hue and is retained. The pipeline, including its parameters and
				probabilities, is not tuned per model.

				\subsubsection{Hyperparameter Tuning and Compute Budget}\label{sec:tuning}

					\textit{Inner split.} Within each outer training fold, one grouped inner split
					divides the data about 80:20 into inner-training and inner-validation sets. It
					uses the same stratified group splitter as the outer folds, so that no
					near-duplicate cluster straddles the two. The inner-validation set alone governs
					every data-dependent choice in that fold:

					\begin{itemize}
						\item when linear probing stops;
						\item early stopping and checkpoint selection in fine-tuning and classifier
							retraining;
						\item the choice of hyperparameters.
					\end{itemize}

					The outer test fold is used for none of them, which avoids the bias
					of tuning and evaluating on the same data
					\citep{cawley2010,raschka2018}.

					The settings are listed in Table~\ref{tab:training}. Hyperparameters are tuned by random search
					\citep{bergstra2012}, and every model receives the same number of search trials, $B$, since unequal
					tuning effort can masquerade as a difference between methods \citep{melis2018}. Phase lengths are not
					searched. In each fold, the best trial on the inner-validation set is that fold's model and is not
					retrained.

					\begin{table}[htbp]
					\vspace{4ex}
					\centering
						\caption{Training and tuning settings, identical for all four models}
						\label{tab:training}
						\singlespacing\footnotesize
						\begin{tabular}{p{4.6cm}p{9.2cm}}
						\hline
						\hline
						Setting & Value \\ \hline
						Phases & Linear probing, fine-tuning, classifier retraining \\
						Inner split & One grouped, stratified 80:20 split of each outer training fold \\
						Search & Random search, $B$ trials per model and fold \\
						Linear-probe learning rate & $10^{-4}$ to $3 \times 10^{-2}$, log-uniform \\
						Fine-tuning learning rate & $10^{-6}$ to $10^{-3}$, log-uniform \\
						Classifier-retraining learning rate & $10^{-4}$ to $3 \times 10^{-2}$, log-uniform \\
						Weight decay & $10^{-6}$ to $10^{-2}$, log-uniform \\
						Optimizer; batch size & AdamW; 32 (fixed) \\
						Early stopping & Inner-validation macro-averaged MAE, patience of 8 epochs \\
						Epoch caps & 30 (linear probing), 50 (fine-tuning), 20 (classifier retraining) \\ \hline\hline
						\end{tabular}
					\vspace{4ex}
					\end{table}

					\textit{Compute budget.} $B$ is fixed from a compute budget of 56 GPU-hours on
					the primary workstation, the lower end of the 56--84 hours available.

					\begin{itemize}
						\item Once step 1 of the sensitivity check (see \secref{sec:sensitivity})
							has measured the cost of one training, $B$ is set to the largest value
							for which the projected total fits within 56 hours with a 10\%
							contingency. The total covers the v1 baseline, the primary comparison, and the final models.
						\item $B$ is written into the study record before any ordinal model is
							trained.
						\item Hours beyond 56 are not used to raise $B$. They go to reruns.
					\end{itemize}

				\subsubsection{Final Models}\label{sec:final}

					Cross-validation estimates the performance of a training procedure. The final
					model of each method is therefore obtained by applying that same procedure once
					to the whole cross-validation pool: a grouped inner split of the pool, the same
					tuning budget $B$, and the three training phases with early stopping. This
					follows the practice of refitting with the selected settings on all available
					training data after model selection \citep{raschka2018}.

					Each method's final model is trained with three seeds. The external test result for a method is the mean over its seeds,
					with their range, so that training variation is separated from external test
					sampling variation. All methods, including the v1 port, are scored on the external test set once and descriptively. The
					external test set is never used to choose among them. 


			\subsection{Evaluation Methods}\label{sec:evaluation}

				\subsubsection{Selection Metric}\label{sec:selmetric}

					Models are selected on macro-averaged MAE, computed per true class on merged
					groups and then averaged across classes \citep{baccianella2009}:
					\[
						\mathrm{MAE}^{M} = \frac{1}{G} \sum_{g} \frac{1}{|g|} \sum_{i:\, y_i \in g}
						\lvert \hat{y}_i - y_i \rvert,
					\]
					where $G$ is the number of merged groups (nine if no merge triggers) and $y_i$
					and $\hat{y}_i$ are the reference and predicted scores on the original 1--9
					scale. The inner mean runs over the images whose reference score falls in group
					$g$. Merging therefore changes only how classes are weighted: a reference 9
					predicted as 8 still counts as an error of one. Macro-averaged mean squared
					error, $\mathrm{MSE}^{M}$, is defined in the same way with squared errors
					\citep{baccianella2009}. Quadratic weighted kappa is deliberately not the selection metric, for the reasons given under \secref{sec:evalmethods}. Because a thin class can move macro-averaged MAE, its stability is measured through the repeated cross-validation and bootstrap protocol rather than assumed.

				\subsubsection{Reported Metrics}\label{sec:metrics}

					Reported alongside the selection metric are the measures in
					Table~\ref{tab:metrics}. Every metric is reported for softmax under both
					its median and argmax readouts, with the median used in the primary
					comparisons. Results are also broken down by source and by device, with
					per-class confidence intervals throughout.

					\begin{table}[htbp]
					\vspace{4ex}
					\centering
						\caption{Metrics reported alongside the selection metric}
						\label{tab:metrics}
						\singlespacing\footnotesize
						\begin{tabular}{p{4.6cm}p{9.2cm}}
						\hline
						\hline
						Metric & Purpose \\ \hline
						Quadratic weighted kappa \citep{cohen1968} & Comparability with the medical
						grading and agricultural literature \\
						Macro-averaged MSE \citep{baccianella2009} & Outlier-sensitive measure for
						Specific Objective 2 \\
						Root mean squared error & Comparability with \citet{shi2023corn} \\
						Plain MAE; within-one-level accuracy & Unweighted summaries of error \\
						Band-crossing error rate; resistance-response confusion matrix & Share of
						predictions in a different response category from the reference (see
						\secref{sec:background}), since macro-averaged MAE weights band-crossing and
						within-band errors equally \\
						Severe-damage recall & Share of images with a reference score of 6 or above,
						the susceptible and highly susceptible categories of \citet{prasanna2018},
						that the model also scores at 6 or above. A severe image scored 5 is a miss,
						because that step crosses into the partially resistant category \\
						Exact accuracy; macro F1 & Continuity with v1, subject to the caveats under
						\secref{sec:baseline} \\
						Full confusion matrices & Per level and per merged group \\ \hline\hline
						\end{tabular}
					\vspace{4ex}
					\end{table}

				\subsubsection{Human Reference}\label{sec:humanref}

					The author's agreement with the expert is reported beside the
					headline figures. This follows a weaker form of the pattern of
					\citet{schaap2022}. Each model and the author are compared under the
					same conditions. Both are scored on exactly the same newly labeled
					images. Both are also checked against the same reference, which is
					the expert's independent score. For expansion images in the
					cross-validation pool, the model's scores are its out-of-fold
					predictions, averaged over the three seeds. For external test images,
					they are the predictions of the final models.

					Two figures are reported side by side for each model. The first is
					the model's macro-averaged MAE against the expert. The second is the
					author's macro-averaged MAE against the expert on the same images.
					Each comes with a 95\% cluster bootstrap interval (see
					\secref{sec:stats}), and quadratic weighted kappa is reported with
					them. The comparison is descriptive, so no margin or verdict is
					applied. If a model's agreement approaches the author's, the model
					rates about as consistently as a person trained on the rubric. This
					does not show that it rates as an expert would (see
					\secref{sec:agreement}). Expansion labeling concentrates on scores 6
					through 9, so the expansion-image result describes the severe range
					only.

			\subsection{Statistical Analysis}\label{sec:statanalysis}

				\subsubsection{Statistical Comparison}\label{sec:stats}

					Because repeated cross-validation folds share training data, a naive paired
					\textit{t}-test across the fifteen runs is anti-conservative
					\citep{dietterich1998}. Model comparisons therefore use the corrected resampled
					\textit{t}-test \citep{nadeau2003}. Bootstrap confidence intervals resample at
					the level of the near-duplicate cluster rather than the image, since image-level
					resampling on data containing near-duplicates yields intervals that are too
					narrow. Because grouped folds vary slightly in size, the test-to-training size
					ratio in the correction is taken as its mean across folds. Bootstrap intervals
					use 2,000 resamples and the percentile method. Clusters are resampled within
					score strata, so that every score is present in every resample, as
					macro-averaged metrics require. 

				\subsubsection{Sensitivity Check Before the Model Comparison}\label{sec:sensitivity}

					Before any ordinal model is trained, a sensitivity power analysis
					\citep{lakens2022} checks whether the design, at its fixed sample size, power,
					and $\alpha$, can detect the smallest effect of interest defined below. Its
					steps, and the responses if it fails, are fixed here.

					\textit{Smallest effect of interest.} The smallest effect of interest is a
					reduction of 0.10 in macro-averaged MAE, for an ordinal model relative to the
					median-readout softmax control. That is one fewer one-step error for every ten
					images in each class. It is set by what would matter in use, not estimated from
					preliminary runs, since power analyses built on pilot effect sizes are biased
					\citep{albers2018}, and it is not revised after any result is seen. The baseline
					run in step 1 is used only for its variance.

					\textit{Primary comparisons.} The check is computed for the three primary
					comparisons: CORAL, CORN, and THOR, each against the median-readout softmax
					control on macro-averaged MAE. Each is assessed at $\alpha' = 0.05/m$ for $m$
					primary comparisons, the strictest step of the Holm procedure that governs the
					tests. Because the variance bound in step 1 applies to any pair of models, the
					check also covers the secondary comparisons of Specific Objective 2.

					\textit{Step 1: measure the variance on the current images.} The v1
					recipe is run on the current 2,084 images under the rebuilt
					protocol, which is arm (c) of the protocol-effect design (see
					\secref{sec:armdesign}) and so needs no extra training. Its fifteen
					per-fold macro-averaged MAE values, under the median readout, give
					the variance from which the smallest effect detectable at 80\%
					power, $\delta$, is computed with the corrected resampled standard
					error \citep{nadeau2003}; the derivation is given in
					Appendix~C.

					\textit{Step 2: project to the planned 2,230 images.} The expansion
					target adds 146 images, all at scores 7--9 (see
					\secref{sec:expansion}), so $\delta$ is projected class by class to
					the planned per-fold test counts (see Appendix~C).
					Once expansion is complete, the detectable effect is recomputed from
					the actual baseline run (arm (d)) and reported with the results,
					whatever it shows.

					\textit{Step 3: decide.} The check passes if the projected $\delta$ is no larger
					than the smallest effect of interest. If it passes, the comparison proceeds as
					planned. If it fails, the following responses are taken in this order, and the
					first that passes is adopted:

					\begin{enumerate}
						\item \textit{Expand further.} The projection is repeated at larger planned
							counts to find the expansion that would pass. This response is taken if
							those labels can be obtained from the expansion pool and local
							collection within the study's timeline.
						\item \textit{Report the ordinal comparison as exploratory.} All four models
							are still trained and reported with their confidence intervals, but no
							significance test is reported as confirmatory, and no conclusion is
							drawn that ordinal training does or does not improve severity scoring.
							Pre-specifying which analyses are confirmatory keeps hypothesis testing
							distinct from hypothesis generation \citep{nosek2018}. Under this
							outcome, the rebuilt protocol, its measured effect (Specific Objective
							4), and the label-consistency findings become the study's primary
							contribution.
					\end{enumerate}

					The check, its inputs, and the response taken are reported in full.

				\subsubsection{Hypothesis Tests}\label{sec:hypotheses}

					Two families of one-sided tests are specified (Table~\ref{tab:hypotheses}).
					Each has familywise $\alpha = 0.05$, controlled by Holm's procedure \citep{holm1979}.

					\begin{itemize}
						\item \textit{Family A} carries Specific Objective 1.
						\item \textit{Family B} carries the two published expectations of Specific
							Objective 2, in the directions reported by \citet{shi2023corn} and
							\citet{fuchs2022}. The direction of the MSE test comes from text-review
							data only; on image data, \citet{shi2023corn} report CORN ahead on both
							MAE and RMSE.
					\end{itemize}

					Each test uses the corrected resampled \textit{t}-test (see \secref{sec:stats}),
					with softmax read out by its median.

					\begin{table}[htbp]
					\vspace{4ex}
					\centering
						\caption{Pre-registered hypothesis tests}
						\label{tab:hypotheses}
						\begin{tabular}{lcll}
						\hline
						\hline
						Family & Objective & One-sided comparison & Metric \\ \hline
						A (primary) & 1 & CORAL better than softmax & Macro-averaged MAE \\
						A (primary) & 1 & CORN better than softmax & Macro-averaged MAE \\
						A (primary) & 1 & THOR better than softmax & Macro-averaged MAE \\
						B (secondary) & 2 & CORN better than CORAL & Macro-averaged MAE \\
						B (secondary) & 2 & CORAL better than CORN & Macro-averaged MSE \\
						B (secondary) & 2 & THOR better than CORAL & Macro-averaged MAE
						\\ \hline\hline
						\end{tabular}
					\vspace{4ex}
					\end{table}

					Each Family A comparison is also given an equivalence test: two one-sided tests
					against bounds of $\pm 0.10$ macro-averaged MAE, the smallest effect of
					interest, using the same corrected standard error. Equivalence is shown when the
					$(1-2\alpha')$ interval of the difference lies entirely within the bounds, with
					$\alpha'$ Holm-adjusted across the family. An equivalence test asks whether an
					effect is smaller than the smallest effect of interest, which is what allows a
					null result to be informative \citep{lakens2017,lakens2018}. Each Family A
					comparison is reported under one of four outcomes:

					\begin{itemize}
						\item significant and not equivalent: the ordinal model is better by a
							meaningful margin;
						\item significant and equivalent: it is better, but by less than the
							smallest effect of interest;
						\item not significant and equivalent: a meaningful improvement is ruled out;
						\item not significant and not equivalent: inconclusive.
					\end{itemize}

					Family B carries no equivalence test, since the smallest effect of interest is
					defined on macro-averaged MAE alone.  All other comparisons in the study are descriptive: they are
					reported with estimates and confidence intervals but no significance test.

		%=====================================================================
		%Chapters IV and V are written after the experiments
		\section{RESULTS AND DISCUSSION}\label{sec:results}

			\textit{To be written after the experiments are complete.}

		\section{SUMMARY AND CONCLUSION}\label{sec:conclusion}

			\textit{To be written after the experiments are complete.}

		%References command
		%Parameter is the references bib file without extension
		\makereferences{references}

		%=====================================================================
		%Appendices
		\section{APPENDICES}\label{sec:appendices}
		\doublespacing

			\subsection{Appendix A. Near-Duplicate Clustering Procedure}\label{app:neardup}

				Near-duplicate clusters (see \secref{sec:leakctrl}) are built in five
				steps:

				\begin{itemize}
					\item \textit{Hash.} A 64-bit DCT perceptual hash is computed after EXIF
						orientation, and the smaller Hamming distance of each image and its
						horizontal mirror is used.
					\item \textit{Calibrate.} The threshold $T$ is the smallest distance that
						captures at least 95\% of synthetic copies of 200 random images (JPEG
						quality 70--95, resized to 50--100\%, cropped by up to 10\%) while
						admitting under 0.1\% of pairs from sources that cannot share a plant,
						such as v1 field and YEESI images.
					\item \textit{Inspect.} Pairs within four bits of $T$ are reviewed with
						labels hidden; $T$ is adjusted once if they show systematic errors,
						then frozen before folds are built.
					\item \textit{Cluster.} Clusters are connected components of pairs at
						distance $\le T$; any cluster over 20 images is checked for chaining
						and split where it joins unrelated images.
					\item \textit{Second pass.} Because hashing misses the same plant
					photographed from a new angle, each image's five nearest neighbors
					by cosine similarity of frozen backbone features are reviewed for a
					random sample of 200 images. A sampled image is a miss if a neighbor
					is confirmed to be the same plant but lies in a different cluster.
					If 5 or more of the 200 (2.5\%) are misses, the review is extended
					to every image. Confirmed same-plant pairs are added to the clusters
					in either case.
				\end{itemize}

				The threshold, calibration distributions, inspection counts, cluster sizes, and the second-pass miss count and additions are reported with the dataset.

			\subsection{Appendix B. Audit of the Inherited Partition}\label{app:partition}

				\setcounter{table}{0}\renewcommand{\thetable}{B.\arabic{table}}

				A file-level audit of the inherited partition, conducted by direct count and
				hashing of the dataset directories rather than from the v1 study's reported
				figures, establishes the distribution in Table~\ref{tab:partition}.

				\begin{table}[htbp]
				\vspace{4ex}
				\centering
					\caption{Inherited train/validation/test partition by severity score}
					\label{tab:partition}
					\begin{tabular}{ccccc}
					\hline
					\hline
					Score & Train & Validation & Test & Total \\ \hline
					1 & 186 & 166 & 33 & 385 \\
					2 & 232 & 45 & 41 & 318 \\
					3 & 434 & 40 & 13 & 487 \\
					4 & 324 & 37 & 28 & 389 \\
					5 & 168 & 25 & 21 & 214 \\
					6 & 127 & 10 & 0 & 137 \\
					7 & 78 & 15 & 5 & 98 \\
					8 & 32 & 15 & 5 & 52 \\
					9 & 4 & 0 & 0 & 4 \\ \hline
					Total & 1,585 & 353 & 146 & 2,084 \\ \hline\hline
					\end{tabular}
				\vspace{4ex}
				\end{table}

				Consequently, the partition has four distinct defects:

				\begin{itemize}
					\item \textit{The validation split is not distributionally representative.}
						Score 1 is 11.7\% of training but 47.0\% of validation; score 3 is
						27.4\% of training but 8.9\% of test.
					\item \textit{Two classes are unevaluable.} Score 6 has 127 training images
						and none in test, so its test recall is undefined. Score 9 appears only
						in training and has never been validated or tested.
					\item \textit{Checkpoint selection was biased.} Because validation results
						guided model selection in v1 \citep{casasola2026}, selection was
						performed against a distribution dominated by score 1---the class with
						the least operational significance, since it denotes no visible damage.
					\item \textit{Training and evaluation distributions diverge.} Scores 5--9
						constitute 25.8\% of the training split but 47.1\% of the 68-image set
						on which the CNN was evaluated \citep{casasola2026}, over-representing
						severe damage by approximately twofold.
				\end{itemize}

				These defects lie in how the data were divided, not in the data or their class
				imbalance, so any comparison computed on this partition measures the partition,
				and it is discarded rather than patched. The v1 study makes no claim that its
				split was stratified; this is an extension of that work, not a critique.

			\subsection{Appendix C. Pre-Registered Analysis Plan: Derivations}\label{app:analysisplan}

				This appendix gives the computations behind steps 1 and 2 of the
				sensitivity check (see \secref{sec:sensitivity}) and the tuning fallback
				of the compute budget (see \secref{sec:tuning}).

				\textit{Step 1: measure the variance on the current images.} The v1 recipe is
				run on the current 2,084 images under the rebuilt protocol, which is arm (c) of
				the protocol-effect design (see \secref{sec:armdesign}) and so needs no extra
				training. Its fifteen per-fold macro-averaged MAE values, under the median
				readout, give a variance estimate $\hat{\sigma}^2$. The variance of a per-fold
				\textit{difference} between two models, unobservable until both are trained, is
				bounded by $2\hat{\sigma}^2$, its value if their errors were uncorrelated. Models
				trained on the same folds make positively correlated errors, so the bound is
				conservative. The corrected standard error of a difference is then
				\[
					\mathrm{SE} = \sqrt{\left(\frac{1}{J} +
					\frac{n_{\text{test}}}{n_{\text{train}}}\right) \cdot 2\hat{\sigma}^2},
				\]
				with $J = 15$ runs and $n_{\text{test}}/n_{\text{train}}$ the mean ratio across
				folds \citep{nadeau2003}, and the smallest detectable effect at 80\% power is
				\[
					\delta = \left(t_{1-\alpha'/2,\,J-1} + t_{0.80,\,J-1}\right) \cdot
					\mathrm{SE}.
				\]
				The baseline's variance, from v1's recipe, stands in for that of the models
				compared, an approximation reported as such.

				\textit{Step 2: project to the planned 2,230 images.} The expansion target adds
				146 images, all at scores 7--9 (see \secref{sec:expansion}). Macro-averaged MAE
				weights every class equally, so its variance is dominated by the classes with
				the fewest test images, which are the classes expansion fills. The projection
				therefore works class by class. Each class's contribution to $\hat{\sigma}^2$ is
				scaled by the ratio of its current to its planned per-fold test count, and SE
				and $\delta$ are recomputed. Score 9, now four images and merged with score 8,
				enters the projection as its own class at its planned count, with the per-image
				error variance measured for score 8. The scaling treats each class's variance as
				arising from the size of its test sample. Variation due to training is also
				reduced by more data, but less predictably. Once
				expansion is complete, the detectable effect is recomputed from the actual
				baseline run (arm (d)) and reported with the results, whatever it shows.

				\textit{Tuning fallback.} If the compute budget gives $B < 5$ for the primary comparison, tuning is instead
					nested in the first seed's five folds only. The configuration chosen
					most often there is reused for the other two seeds, and the resulting
					optimistic bias is reported.

			\subsection{Appendix D. FAWCheck v1 Recipe and Its PyTorch Port}\label{app:recipe}

				\setcounter{table}{0}\renewcommand{\thetable}{D.\arabic{table}}

				Table~\ref{tab:recipe} lists each setting of the v1 recipe and its
				PyTorch counterpart (see \secref{sec:recipe}). Some settings are not
				recorded in the v1 model file. Their values are stated defaults, and
				the table marks them as such.

				\begin{table}[htbp]
				\vspace{4ex}
				\centering
					\caption{FAWCheck v1 training recipe, as recorded in the v1 model file, and
					its PyTorch port}
					\label{tab:recipe}
					\singlespacing\footnotesize
					\begin{tabular}{p{2.6cm}p{5.6cm}p{5.6cm}}
					\hline
					\hline
					Setting & v1 (Keras 3.13.2) & PyTorch port \\ \hline
					Backbone & MobileNetV3Small, ImageNet weights, without its classification
					top & torchvision MobileNetV3-Small, ImageNet weights, classifier removed
					\\ \hline
					Input & $224 \times 224$ RGB, 0--255, rescaled to $[-1, 1]$ inside the
					network & $224 \times 224$ RGB, scaled to $[0, 1]$ and normalized with
					ImageNet mean and standard deviation \\ \hline
					Trainable layers & Last three inverted-residual blocks and final $1 \times
					1$ convolution; all earlier layers frozen & The same blocks (torchvision
					\texttt{features[9]}--\texttt{features[12]});
					\texttt{features[0]}--\texttt{features[8]} frozen \\ \hline
					Batch normalization & Frozen layers keep pretrained statistics; momentum
					0.999, $\varepsilon = 0.001$ & Frozen blocks in evaluation mode; trainable
					blocks update statistics, momentum 0.001 (PyTorch convention), $\varepsilon
					= 0.001$ \\ \hline
					Output layer & Global average pooling $\rightarrow$ dropout 0.3
					$\rightarrow$ dense, 9 units, softmax; Glorot-uniform initialization &
					Global average pooling $\rightarrow$ dropout 0.3 $\rightarrow$ linear, 9
					units; Xavier-uniform initialization; softmax at readout \\ \hline
					Loss & Sparse categorical cross-entropy & Cross-entropy \\ \hline
					Optimizer & Adam, learning rate $1 \times 10^{-5}$, $\beta_1 = 0.9$,
					$\beta_2 = 0.999$, $\varepsilon = 1 \times 10^{-7}$ & Same \\ \hline
					Augmentation & Horizontal flip; rotation up to $\pm 36^{\circ}$; zoom up to
					$\pm 10\%$, aspect ratio preserved; contrast factor between 1.0 and 1.1 per
					image, about each channel's mean; reflect fill, bilinear interpolation &
					Same transformations and parameters, applied in the same order after
					resizing to $224 \times 224$ \\ \hline
					Readout & Argmax & Argmax; median also reported \\ \hline
					Class weighting & Not recorded & None (default) \\ \hline
					Batch size & Not recorded & 32 (default) \\ \hline
					Epochs and stopping & Not recorded; checkpoint selection guided by
					validation results \citep{casasola2026} & Up to 50 epochs; early stopping on
					validation loss with patience 5, best weights restored (default) \\ \hline
					Output-layer-only phase before fine-tuning & Not recorded & None (default)
					\\ \hline\hline
					\end{tabular}
				\vspace{4ex}
				\end{table}

				v1's contrast augmentation only ever increases contrast: Keras 3 bounds
				the lower contrast factor at zero, so a factor of 0.1 yields multipliers
				between 1.0 and 1.1. The port reproduces this rather than correcting it.

	\end{mainmatter}
\end{document}
